\documentclass{subfiles}
\begin{document}
    Let \(\A\) be an algorithm and let \(f(n)\) be the function that describes how 
        \(\A\) behaves on an input of size \(n\); by simply analyzing \(f\),
        we can deduce how efficient the whole algorithm is. 
        In this context three scenarios must be considered, namely,
        how does \(\A\) behaves in the best, the average and worst case.
        
    Each of the three scenarios corresponds to a different notation that takes in account 
        just \(f\). These are the \(\Omega, \Theta \text{ and } \mathcal{O}\) notations.

    Formally speaking, given \(g(n)\) a function, we define
    \[\begin{gathered}
        \omegaOf{g(n)} = \set*{f(n)}[\exists c, n_{0} > 0 : \forall n \ge n_{0} \ f(n) \ge cg(n)]\\
        \bigO{g(n)} = \set*{f(n)}[\exists c, n_{0} > 0 : \forall n \ge n_{0} \ f(n) \le cg(n)] \\ 
        \thetaOf{g(n)} = \set*{f(n)}[\exists c_{1}, c_{2}, n_{0} > 0 : \forall n \ge n_{0} \ c_{1}g(n) \le f(n) \le c_{2}g(n)].
    \end{gathered}\]
    Despite given in set form, we usually write \(f(n) = \bigO{g(n)}\) (and the same hold for the other notations),
        instead of \(f(n) \in \bigO{g(n)}\). That is, asymptotic notations measure how 
        functions behave without caring of multiplicative constants.
\end{document}
